%
% this file is encoded in utf-8
% v2.0 (Apr. 5, 2009)
% do not change the content of this file
% unless the thesis layout rule is changed
% 無須修改本檔內容，除非校方修改了
% 封面、書名頁、中文摘要、英文摘要、誌謝、目錄、表目錄、圖目錄、符號說明
% 等頁之格式

% make the line spacing in effect
\renewcommand{\baselinestretch}{\mybaselinestretch}
\large % it needs a font size changing command to be effective

% default variables definitions
% 此處只是預設值，不需更改此處
% 請更改 my_names.tex 內容
\newcommand\cTitle{論文題目}
\newcommand\eTitle{MY THESIS TITLE}
\newcommand\myCname{....}
\newcommand\myEname{Aron Wang}
\newcommand\myStudentID{M9315048}    % add student id by saiba 
\newcommand\advisorCnameA{南宮明博士}
\newcommand\advisorEnameA{Dr.~Ming Nangong}
\newcommand\advisorCnameB{李斯坦博士}
\newcommand\advisorEnameB{Dr.~Stein Lee}
\newcommand\advisorCnameC{徐　石博士}
\newcommand\advisorEnameC{Dr.~Sean~Hsu}
\newcommand\univCname{元智大學}
\newcommand\univEname{Yuan Ze University}
\newcommand\deptCname{光電工程研究所}
\newcommand\fulldeptEname{Graduate School of Electro-Optical Engineering}
\newcommand\deptEname{Electro-optical Engineering}
\newcommand\collEname{College of Engineering}
\newcommand\degreeCname{碩士}
\newcommand\degreeEname{Master of Science}
\newcommand\cYear{九十四}
\newcommand\cMonth{六}
\newcommand\cDay{六}
\newcommand\eYear{2006}
\newcommand\eMonth{June}
\newcommand\ePlace{Chungli, Taoyuan, Taiwan}


 % user's names; to replace those default variable definitions
%
% this file is encoded in utf-8
% v2.0 (Apr. 5, 2009)
% 填入你的論文題目、姓名等資料
% 如果題目內有必須以數學模式表示的符號，請用 \mbox{} 包住數學模式，如下範例
% 如果中文名字是單名，與姓氏之間建議以全形空白填入，如下範例
% 英文名字中的稱謂，如 Prof. 以及 Dr.，其句點之後請以不斷行空白~代替一般空白，如下範例
% 如果你的指導教授沒有如預設的三位這麼多，則請把相對應的多餘教授的中文、英文名
%    的定義以空的大括號表示
%    如，\renewcommand\advisorCnameB{}
%          \renewcommand\advisorEnameB{}
%          \renewcommand\advisorCnameC{}
%          \renewcommand\advisorEnameC{}

% 論文題目 (中文)
\renewcommand\cTitle{%
採用深度強化學習結合時間序列神經網路預測模型的自適性交通信號系統設計
}

% 論文題目 (英文)
\renewcommand\eTitle{%
Design of Adaptive Traffic Light Signaling System Using Deep Reinforcement Learning combined with Time Series Neural Network Forecasting Model}

% 我的姓名 (中文)
\renewcommand\myCname{吳宜真}

% 我的姓名 (英文)
\renewcommand\myEname{Yi-Chen Wu}

% 我的學號 (英文)
\renewcommand\myStudentID{M10815045}

% 指導教授A的姓名 (中文)
\renewcommand\advisorCnameA{馮輝文 \quad 博士}

% 指導教授A的姓名 (英文)
\renewcommand\advisorEnameA{Dr. Huei-Wen Ferng}

% 指導教授B的姓名 (中文)
\renewcommand\advisorCnameB{}

% 指導教授B的姓名 (英文)
\renewcommand\advisorEnameB{}

% 指導教授C的姓名 (中文)
\renewcommand\advisorCnameC{}

% 指導教授C的姓名 (英文)
\renewcommand\advisorEnameC{}

% 校名 (中文)
\renewcommand\univCname{國立臺灣科技大學}

% 校名 (英文)
\renewcommand\univEname{National Taiwan University of Science and Technology}

% 系所名 (中文)
\renewcommand\deptCname{資訊工程系}

% 系所全名 (英文)
\renewcommand\fulldeptEname{Department of Computer Science and Information Engineering}

% 系所短名 (英文, 用於書名頁學位名領域)
%\renewcommand\deptEname{Electro-Optical Engineering}

% 學院英文名 (如無，則以空的大括號表示)
\renewcommand\collEname{College of Electrical Engineering and Computer Science}

% 學位名 (中文)
\renewcommand\degreeCname{碩士}

% 學位名 (英文)
\renewcommand\degreeEname{Master of Science}

% 口試年份 (中文、民國)
\renewcommand\cYear{一百一十一}

% 口試月份 (中文)
\renewcommand\cMonth{十} 

% 口試日期 (中文)
\renewcommand\cDay{十九}

% 口試年份 (阿拉伯數字、西元)
\renewcommand\eYear{2022} 

% 口試月份 (英文)
\renewcommand\eMonth{October}

% 學校所在地 (英文)
%\renewcommand\ePlace{Taipei City, Taiwan}

%畢業級別；用於書背列印；若無此需要可忽略
\newcommand\GraduationClass{103}

%%%%%%%%%%%%%%%%%%%%%%

% 使用 hyperref 在 pdf 簡介欄裡填入相關資料
\ifx\hypersetup\undefined
	\relax  % do nothing
\else
	\hypersetup{
	pdftitle=\cTitle,
	pdfauthor=\myCname}
\fi
	

\newcommand\itsempty{}
%%%%%%%%%%%%%%%%%%%%%%%%%%%%%%%
%       NTUST cover 封面
%%%%%%%%%%%%%%%%%%%%%%%%%%%%%%%
%
\begin{titlepage}
% no page number
% next page will be page 1

% aligned to the center of the page
\begin{center}
% font size (relative to 12 pt):
% \large (14pt) < \Large (18pt) < \LARGE (20pt) < \huge (24pt)< \Huge (24 pt)
%

\begin{figure}[htbp]
	\begin{minipage}[b]{5cm} 
		\raggedright
		\includegraphics[width=1.5in]{frontpages/ntust_logo}
		\label{fig:ntust_logo}
	\end{minipage}% 
	\begin{minipage}[b]{0.5\textwidth} 
	\centering
	\makebox[3cm][c]{\Huge{\univCname}}\\  %顯示中文校名
	\vspace{0.5cm}
	\makebox[3cm][c]{\Huge{\deptCname}}\\ %顯示中文系所名
	\vspace{0.5cm}
	\end{minipage}%
\\ 
\rule{16cm}{3pt}
\end{figure}
%\hfill

\vspace{1cm}
\makebox[6cm][s]{\textbf{\Huge{\degreeCname 論文}}}\\ %顯示論文種類 (中文)
\vspace{1cm}
%
% Set the line spacing to single for the titles (to compress the lines)
\renewcommand{\baselinestretch}{1}   %行距 1 倍
%\large % it needs a font size changing command to be effective
\Large{\cTitle}\\  % 中文題目
%
\vspace{1cm}
%
\Large{\eTitle}\\ %英文題目
\vspace{5cm}
% \makebox is a text box with specified width;
% option s: stretch
% use \makebox to make sure
% 「研究生：」 與「指導教授：」occupy the same width
%
% remark by saiba
% 由於全型符號在xelatex中無法正確的左右分散，所以將中文與符號部分分開處理

\hspace{4cm} \makebox[3cm][s]{\Large{研究生}}\makebox[.5cm][s]{\Large{：}}
\Large{\myCname}  % 顯示作者中文名
\hfill \makebox[1cm][s]{}\\
%
\vspace{0.3cm}
\hspace{4cm} \makebox[3cm][s]{\Large{學號}}\makebox[.5cm][s]{\Large{：}}
\Large{\myStudentID}  %顯示指導教授A中文名
\hfill \makebox[1cm][s]{}\\
%
\vspace{1cm}
\hspace{4cm} \makebox[3cm][s]{\Large{指導教授}}\makebox[.5cm][s]{\Large{：}}
\Large{\advisorCnameA}  %顯示指導教授A中文名
\hfill \makebox[1cm][s]{}\\
%
% 判斷是否有共同指導的教授 B
\ifx \advisorCnameB  \itsempty
\relax % 沒有 B 教授，所以不佔版面，不印任何空白
\else
% 共同指導的教授 B
\hspace{4.5cm} \makebox[3cm][s]{}
\Large{\advisorCnameB}  %顯示指導教授B中文名
\hfill \makebox[1cm][s]{}\\
\fi
%
% 判斷是否有共同指導的教授 C
\ifx \advisorCnameC  \itsempty
\relax % 沒有 C 教授，所以不佔版面，不印任何空白
\else
% 共同指導的教授 C
\hspace{4.5cm} \makebox[3cm][s]{}
\Large{\advisorCnameC}  %顯示指導教授C中文名
\hfill \makebox[1cm][s]{}\\
\fi
%
\vfill
\makebox[10cm][s]{\Large{中華民國\cYear 年\cMonth 月\cDay 日}}%
%
\end{center}
% Resume the line spacing to the desired setting
\renewcommand{\baselinestretch}{\mybaselinestretch}   %恢復原設定
% it needs a font size changing command to be effective
% restore the font size to normal
\normalsize
\end{titlepage}
%%%%%%%%%%%%%%

%% 從摘要到本文之前的部份以小寫羅馬數字印頁碼
% 但是從「書名頁」(但不印頁碼) 就開始計算
\setcounter{page}{1}
\pagenumbering{roman} 
%\pagenumbering{arabic}



%%%%%%%%%%%%%%%%%%%%%%%%%%%%%%%%
%%       書名頁 
%%%%%%%%%%%%%%%%%%%%%%%%%%%%%%%%
%%
%\newpage
%
%% 判斷是否要浮水印？
%\ifx\mywatermark\undefined 
%  \thispagestyle{empty}  % 無頁碼、無 header (無浮水印)
%\else
%  \thispagestyle{EmptyWaterMarkPage} % 無頁碼、有浮水印
%\fi
%
%%no page number
%% create an entry in table of contents for 書名頁
%\phantomsection % for hyperref to register this
%\addcontentsline{toc}{chapter}{\nameInnerCover}
%
%
%% aligned to the center of the page
%\begin{center}
%% font size (relative to 12 pt):
%% \large (14pt) < \Large (18pt) < \LARGE (20pt) < \huge (24pt)< \Huge (24 pt)
%% Set the line spacing to single for the titles (to compress the lines)
%\renewcommand{\baselinestretch}{1}   %行距 1 倍
%% it needs a font size changing command to be effective
%%中文題目
%\Large{\cTitle}\\ %%%%%
%\vspace{1cm}
%% 英文題目
%\Large{\eTitle}\\ %%%%%
%%\vspace{1cm}
%\vfill
%% \makebox is a text box with specified width;
%% option s: stretch
%% use \makebox to make sure
%% 「研究生：」 與「指導教授：」occupy the same width
%\makebox[3cm][s]{\large{研 究 生：}}
%\makebox[3cm][l]{\large{\myCname}} %%%%%
%\hfill
%\makebox[2cm][s]{\large{Student: }}
%\makebox[5cm][l]{\large{\myEname}}\\ %%%%%
%%
%%\vspace{1cm}
%%
%\makebox[3cm][s]{\large{指導教授：}}
%\makebox[3cm][l]{\large{\advisorCnameA}} %%%%%
%\hfill
%\makebox[2cm][s]{\large{Advisor: }}
%\makebox[5cm][l]{\large{\advisorEnameA}}\\ %%%%%
%%
%% 判斷是否有共同指導的教授 B
%\ifx \advisorCnameB  \itsempty
%\relax % 沒有 B 教授，所以不佔版面，不印任何空白
%\else
%%共同指導的教授B
%\makebox[3cm][s]{}
%\makebox[3cm][l]{\large{\advisorCnameB}} %%%%%
%\hfill
%\makebox[2cm][s]{}
%\makebox[5cm][l]{\large{\advisorEnameB}}\\ %%%%%
%\fi
%%
%% 判斷是否有共同指導的教授 C
%\ifx \advisorCnameC  \itsempty
%\relax % 沒有 C 教授，所以不佔版面，不印任何空白
%\else
%%共同指導的教授C
%\makebox[3cm][s]{}
%\makebox[3cm][l]{\large{\advisorCnameC}} %%%%%
%\hfill
%\makebox[2cm][s]{}
%\makebox[5cm][l]{\large{\advisorEnameC}}\\ %%%%%
%\fi
%%
%% Resume the line spacing to the desired setting
%\renewcommand{\baselinestretch}{\mybaselinestretch}   %恢復原設定
%\large %it needs a font size changing command to be effective
%%
%\vfill
%\makebox[4cm][s]{\large{\univCname}}\\% 校名
%\makebox[6cm][s]{\large{\deptCname}}\\% 系所名
%\makebox[3cm][s]{\large{\degreeCname 論文}}\\% 學位名
%%
%%\vspace{1cm}
%\vfill
%\large{A Thesis}\\%
%\large{Submitted to }%
%%
%\large{\fulldeptEname}\\%系所全名 (英文)
%%
%%
%\ifx \collEname  \itsempty
%\relax % 沒有學院名 (英文)，所以不佔版面，不印任何空白
%\else
%% 有學院名 (英文)
%\large{\collEname}\\% 學院名 (英文)
%\fi
%%
%\large{\univEname}\\%校名 (英文)
%%
%\large{in Partial Fulfillment of the Requirements}\\
%%
%\large{for the Degree of}\\
%%
%\large{\degreeEname}\\%學位名(英文)
%%
%\large{in}\\
%%
%\large{\deptEname}\\%系所短名(英文；表明學位領域)
%%
%\large{\eMonth\ \eYear}\\%月、年 (英文)
%%
%\large{\ePlace}% 學校所在地 (英文)
%\vfill
%\large{中華民國}%
%\large{\cYear}% %%%%%
%\large{年}%
%\large{\cMonth}% %%%%%
%\large{月}\\
%\end{center}
%% restore the font size to normal
%\normalsize
%\clearpage
%%%%%%%%%%%%%%%%%%%%%%%%%%%%%%%%
%%       論文口試委員審定書 (計頁碼，但不印頁碼) 
%%%%%%%%%%%%%%%%%%%%%%%%%%%%%%%%
%%
%% insert the printed standard form when the thesis is ready to bind
%% 在口試完成後，再將已簽名的審定書放入以便裝訂
%% create an entry in table of contents for 審定書
%% 目前送出空白頁
%\newpage%
%{\thispagestyle{empty}%
%\phantomsection % for hyperref to register this
%\addcontentsline{toc}{chapter}{\nameCommitteeForm}%
%\mbox{}\clearpage}

%%%%%%%%%%%%%%%%%%%%%%%%%%%%%%%
%       指導教授推薦書 (計頁碼，但不印頁碼)
%%%%%%%%%%%%%%%%%%%%%%%%%%%%%%%
%
% insert the printed standard form when the thesis is ready to bind
% 在口試完成後，再將已簽名的推薦書放入以便裝訂
% create an entry in table of contents for 推薦書
% 目前送出空白頁
\newpage%
{\thispagestyle{empty}%
\phantomsection % for hyperref to register this
\addcontentsline{toc}{chapter}{\nameRecommForm}%
\mbox{}\clearpage}


%%%%%%%%%%%%%%%%%%%%%%%%%%%%%%%
%       考試委員審定書 (計頁碼，但不印頁碼) 
%%%%%%%%%%%%%%%%%%%%%%%%%%%%%%%
%
% insert the printed standard form when the thesis is ready to bind
% 在口試完成後，再將已簽名的審定書放入以便裝訂
% create an entry in table of contents for 審定書
% 目前送出空白頁
\newpage%
{\thispagestyle{empty}%
\phantomsection % for hyperref to register this
\addcontentsline{toc}{chapter}{\nameCommitteeForm}%
\mbox{}\clearpage}

%%%%%%%%%%%%%%%%%%%%%%%%%%%%%%%
%       中文摘要 
%%%%%%%%%%%%%%%%%%%%%%%%%%%%%%%
%
\newpage
\thispagestyle{plain}  % 無 header，但在浮水印模式下會有浮水印
% create an entry in table of contents for 中文摘要
\phantomsection % for hyperref to register this
\addcontentsline{toc}{chapter}{\nameCabstract}

% aligned to the center of the page
\begin{center}
% font size (relative to 12 pt):
% \large (14pt) < \Large (18pt) < \LARGE (20pt) < \huge (24pt)< \Huge (24 pt)
% Set the line spacing to single for the names (to compress the lines)
\renewcommand{\baselinestretch}{1}   %行距 1 倍
% it needs a font size changing command to be effective
\large{\cTitle}\\  %中文題目
\vspace{0.83cm}
% \makebox is a text box with specified width;
% option s: stretch
% use \makebox to make sure
% each text field occupies the same width
\makebox[1.5cm][s]{\large{學生：}}
\makebox[3cm][l]{\large{\myCname}} %學生中文姓名
\hfill
%
\makebox[3cm][s]{\large{指導教授：}}
\makebox[3cm][l]{\large{\advisorCnameA}} \\ %教授A中文姓名
%
% 判斷是否有共同指導的教授 B
\ifx \advisorCnameB  \itsempty
\relax % 沒有 B 教授，所以不佔版面，不印任何空白
\else
%共同指導的教授B
\makebox[1.5cm][s]{}
\makebox[3cm][l]{} %%%%%
\hfill
\makebox[3cm][s]{}
\makebox[3cm][l]{\large{\advisorCnameB}}\\ %教授B中文姓名
\fi
%
% 判斷是否有共同指導的教授 C
\ifx \advisorCnameC  \itsempty
\relax % 沒有 C 教授，所以不佔版面，不印任何空白
\else
%共同指導的教授C
\makebox[1.5cm][s]{}
\makebox[3cm][l]{} %%%%%
\hfill
\makebox[3cm][s]{}
\makebox[3cm][l]{\large{\advisorCnameC}}\\ %教授C中文姓名
\fi
%
\vspace{0.42cm}
%
\large{\univCname}\large{\deptCname}\\ %校名系所名
\vspace{0.83cm}
%\vfill
\makebox[2.5cm][s]{\large{摘要}}\\
\end{center}
% Resume the line spacing to the desired setting
\renewcommand{\baselinestretch}{\mybaselinestretch}   %恢復原設定
%it needs a font size changing command to be effective
% restore the font size to normal
\normalsize
%%%%%%%%%%%%%

由於目前長期演進上行傳輸使用的經典排程演算法並未顧及使用者裝置的延遲狀況且其公平性不佳，導致部分使用者裝置得不到資源以傳輸資料且封包嚴重遺失，再者，資源分配方式的效率亦不佳，因此，本碩士論文著眼於設計一排程演算法來消弭封包大量遺失的狀況並提升使用者裝置的公平性。為能改善封包遺失率，我們的排程演算法不以使用者裝置通道品質為基礎，而是加入延遲預算做考量，選擇使用者裝置中平均佇列頭端封包剩餘延遲預算最少的使用者裝置為優先對象，使其能優先獲得較佳的資源以進行資料傳輸，降低使用者裝置的封包遺失率。除此之外，透過參考使用者裝置平均配置過的資源數量，讓取得過較少資源的使用者裝置也能有機會取得資源，進而提升整體的公平性，另外，我們的排程演算法在連續分配上設置分配數量限制，能更有效地利用資源，進一步提升系統的吞吐量。最後，在固定服務流數量與隨機服務流數量安排下，模擬數據印證本碩士論文提出之排程演算法均能達到預期的效果，與其他相關的排程演算法相比，除了能有效降低封包的遺失率，更可提升整體使用者裝置彼此間的公平性及系統的吞吐量。
\\
\\
關鍵字：長期演進，上行排程演算法，封包延遲，延遲預算，封包遺失率。
%我們在排程演算法的設計上加入延遲預算做為主要考量，優先分配資源給封包延遲時間逼近延遲預算的使用者，讓排程器能優先選取較迫切需要得到資源的使用者，使其能優先獲得較佳的資源以進行資料傳輸，而其他經典的排程演算法在沒有考量延遲預算的情況下，通常以較佳通道品質為主要考量的資源分配方式，容易使通道品質較差的使用者得不到資源，導致封包因超過其最大延遲時間而被丟棄，造成封包遺失率上升。除此之外，為了在連續分配上能更有效的利用資源，進一步提升系統的吞吐量，我們的排程演算法能(透過什麼原因/方式/作法)讓通道品質較差的使用者也能有機會取得資源，進而提升整體的公平性。

%%%%%%%%%%%%%%%%%%%%%%%%%%%%%%%
%       英文摘要 
%%%%%%%%%%%%%%%%%%%%%%%%%%%%%%%
%
\newpage
\thispagestyle{plain}  % 無 header，但在浮水印模式下會有浮水印

% create an entry in table of contents for 英文摘要
\phantomsection % for hyperref to register this
\addcontentsline{toc}{chapter}{\nameEabstract}

% aligned to the center of the page
\begin{center}
% font size:
% \large (14pt) < \Large (18pt) < \LARGE (20pt) < \huge (24pt)< \Huge (24 pt)
% Set the line spacing to single for the names (to compress the lines)
\renewcommand{\baselinestretch}{1}   %行距 1 倍
%\large % it needs a font size changing command to be effective
\large{\eTitle}\\  %英文題目
\vspace{0.83cm}
% \makebox is a text box with specified width;
% option s: stretch
% use \makebox to make sure
% each text field occupies the same width
\makebox[2cm][s]{\large{Student: }}
\makebox[5cm][l]{\large{\myEname}} %學生英文姓名
\hfill
%
\makebox[2cm][s]{\large{Advisor: }}
\makebox[5cm][l]{\large{\advisorEnameA}} \\ %教授A英文姓名
%
% 判斷是否有共同指導的教授 B
\ifx \advisorCnameB  \itsempty
\relax % 沒有 B 教授，所以不佔版面，不印任何空白
\else
%共同指導的教授B
\makebox[2cm][s]{}
\makebox[5cm][l]{} %%%%%
\hfill
\makebox[2cm][s]{}
\makebox[5cm][l]{\large{\advisorEnameB}}\\ %教授B英文姓名
\fi
%
% 判斷是否有共同指導的教授 C
\ifx \advisorCnameC  \itsempty
\relax % 沒有 C 教授，所以不佔版面，不印任何空白
\else
%共同指導的教授C
\makebox[2cm][s]{}
\makebox[5cm][l]{} %%%%%
\hfill
\makebox[2cm][s]{}
\makebox[5cm][l]{\large{\advisorEnameC}}\\ %教授C英文姓名
\fi
%
\vspace{0.42cm}
\large{Submitted to }\\\large{\fulldeptEname }\\  %英文系所全名
%
\ifx \collEname  \itsempty
\relax % 如果沒有學院名 (英文)，則不佔版面，不印任何空白
\else
% 有學院名 (英文)
\large{\collEname }\\% 學院名 (英文)
\fi
%
\large{\univEname}\\  %英文校名
\vspace{0.83cm}
%\vfill
%
\large{ABSTRACT}\\
%\vspace{0.5cm}
\end{center}
% Resume the line spacing the desired setting
\renewcommand{\baselinestretch}{\mybaselinestretch}   %恢復原設定
%\large %it needs a font size changing command to be effective
% restore the font size to normal
\normalsize
%%%%%%%%%%%%%

Nowadays, classic scheduling algorithms used in LTE uplink fail to take care of delay conditions and fairness among user equipments(UEs), resulting in severe packet losses and poor fairness for some UEs. Meanwhile, the corresponding resource allocation is not efficient as well. Therefore, we aim at designing a scheduling algorithm to avoid packet losses and enhance fairness. Targeting towards the aforementioned goals, the delay budget rather than the channel condition of UEs serves as the key factor when designing our algorithm. Our algorithm employs the principle that the UE with the least delay budget left in queue will be scheduled preferentially first. By doing so, it is allowed to get good-quality resources to transmit data so that the problem of packet losses can be alleviated. Furthermore, the amount of allocated resources is taken into account. UEs with fewer allocated resources still have chances to be served to improve fairness. In addition, our proposed algorithm poses a constraint in sequential resource allocation so that resources can be utilized efficiently and the overall throughput can be improved. Finally, our simulation results support that our proposed algorithm performs well as expected either in a scenario with a fixed number of service flows or in a scenario with a random number of service flows. Compared with the other related algorithms, our proposed algorithm has fewer packet losses as well as better system throughput.
\\
\\
Keywords：LTE, Uplink Scheduling, Packet Delay, Delay Budget, Packet Loss Ratio.
%having packets with shorter time left in queue to be scheduled and sets constraint in sequential allocation method.
%Under circumstances that reliability and penetration of intellegent portable devices rises with technology in modern society, efficient scheduling with limited resources plays an important role.
% To maximize system throughput, algorithms widely used in uplink regard users with good channel conditions as primary resource allocation objects. Those with poor channel conditions are lack of resources and drop packets frequently that results in high packet loss ratio, poor throughput and fairness. The method of sequential resource block allocation effects the throughput as well.

%%%%%%%%%%%%%%%%%%%%%%%%%%%%%%%
%       誌謝 
%%%%%%%%%%%%%%%%%%%%%%%%%%%%%%%
%
% Acknowledgment
\newpage
\chapter*{\protect\makebox[5cm][s]{\nameAckn}} %\makebox{} is fragile; need protect
\phantomsection % for hyperref to register this
\addcontentsline{toc}{chapter}{\nameAckn}
首先，我要向我的指導教授\underline{馮輝文} 教授致上誠摯的感謝，在這兩年多以來的碩士學業中，不只是在學術知識上給予指導，也關心學生們的身心狀況，同時期望我們不只在面對學術研究時能有嚴謹且實事求是的態度，在面對生活與做人處事上更能認真負責，抱持積極正向的心境走過遭遇的困境。更感謝教授在我輾轉找尋研究方向時，總能耐心的與我討論並指出正確的思考方向，最後能順利找到研究方向並深入探討，同時也在論文的撰寫上給予了我許多建議與方向，兩年多來的研究生活著實紮實且獲益匪淺，非常感謝教授對於我的教誨與付出的心力。

同時也要感謝口試委員\underline{林嘉慶}教授、\underline{謝宏昀}教授以及\underline{黎碧煌}教授能在繁忙的事務中抽出時間，同時感謝教授們在報告後能以不同的思維角度，給予我諸多專業且具啟發性的意見與只導，使得研究的內容能更加充實，本論文得以更臻完善。

此外，也要感謝無線通訊暨網路工程實驗室的所有成員，感謝\underline{許雅淑}學姊、\underline{雷正}學長、\underline{趙清松}學長、\underline{蔡亞宸}學姊、\underline{黃雅喻}學姊、\underline{王柏崴}學長與\underline{楊育維}學長對我的諸多照顧，總是給予許多經驗的傳授及學術知識的指導，讓我獲益良多。同時也感謝\underline{陳佳瑩}、\underline{王騰輝}與\underline{蘇毓傑}陪同一起修課，並在研究的過程中彼此給予鼓勵與安慰。也謝謝\underline{李姵瑩}、\underline{蔡雅如}、\underline{洪家琪}、\underline{劉羽軒}、\underline{曾英祖}、\underline{馬維均}、\underline{黃光緯}、\underline{鄭丁元}等學弟妹們幫忙分擔處理實驗室的各種事務，並帶來許多歡笑與輕鬆舒適的氛圍。同時也謝謝\underline{常家銘}、\underline{陳信安}、\underline{戴邵軒}、\underline{黃柏崴}、\underline{夏致群}、\underline{劉澤}、\underline{劉冠甫}等大學時的同學及學長們，在研究生生活上的各種相互扶持與幫助。

最後，要感謝我的家人對我的照顧，不時給予我精神的鼓勵與物質上的支持，讓我能在求學過程中能專心進行研究，也謝謝女友\underline{琳達}一直以來的陪伴，互相分享生活上發生的大小事，即使在外地工作仍能感受心中的溫暖。感謝所有幫助過我的人，若非有你們，我無法獨自完成今日的成果。願家人和所有愛我與我愛的人，能一同分享碩士研究與論文完成的喜悅。


\begin{flushright}
李承洋\quad 謹誌於 2016/12/19 臺灣科大 無線通訊暨網路工程實驗室RB306-2
\end{flushright}
%\raggedleft 

%%%%%%%%%%%%%%%%%%%%%%%%%%%%%%%
%       目錄 
%%%%%%%%%%%%%%%%%%%%%%%%%%%%%%%
%
% Table of contents
\newpage
\renewcommand{\contentsname}{\protect\makebox[5cm][s]{\nameToc}}
%\makebox{} is fragile; need protect
\phantomsection % for hyperref to register this
\addcontentsline{toc}{chapter}{\nameToc}
\tableofcontents

%%%%%%%%%%%%%%%%%%%%%%%%%%%%%%%
%       表目錄 
%%%%%%%%%%%%%%%%%%%%%%%%%%%%%%%
%
% List of Tables
\newpage
\renewcommand{\listtablename}{\protect\makebox[5cm][s]{\nameLot}}
%\makebox{} is fragile; need protect
\phantomsection % for hyperref to register this
\addcontentsline{toc}{chapter}{\nameLot}
\listoftables

%%%%%%%%%%%%%%%%%%%%%%%%%%%%%%%
%       圖目錄 
%%%%%%%%%%%%%%%%%%%%%%%%%%%%%%%
%
% List of Figures
\newpage
\renewcommand{\listfigurename}{\protect\makebox[5cm][s]{\nameTof}}
%\makebox{} is fragile; need protect
\phantomsection % for hyperref to register this
\addcontentsline{toc}{chapter}{\nameTof}
\listoffigures

%%%%%%%%%%%%%%%%%%%%%%%%%%%%%%%
%       符號說明 
%%%%%%%%%%%%%%%%%%%%%%%%%%%%%%%
%
% Symbol list
% define new environment, based on standard description environment
% adapted from p.60~64, <<The LaTeX Companion>>, 1994, ISBN 0-201-54199-8
\newcommand{\SymEntryLabel}[1]%
  {\makebox[3cm][l]{#1}}
%
\newenvironment{SymEntry}
   {\begin{list}{}%
       {\renewcommand{\makelabel}{\SymEntryLabel}%
        \setlength{\labelwidth}{3cm}%
        \setlength{\leftmargin}{\labelwidth}%
        }%
   }%
   {\end{list}}
%%
%\newpage
%\chapter*{\protect\makebox[5cm][s]{\nameSlist}} %\makebox{} is fragile; need protect
%\phantomsection % for hyperref to register this
%\addcontentsline{toc}{chapter}{\nameSlist}
%%
% this file is encoded in utf-8
% v2.0 (Apr. 5, 2009)
%  各符號以 \item[] 包住，然後接著寫說明
% 如果符號是數學符號，應以數學模式$$表示，以取得正確的字體
% 如果符號本身帶有方括號，則此符號可以用大括號 {} 包住保護
\begin{SymEntry}

\item[OLED]
Organic Light Emitting Diode

\item[$E$]
energy

\item[$e$]
the absolute value of the electron charge, $1.60\times10^{-19}\,\text{C}$
 
\item[$\mathscr{E}$]
electric field strength (V/cm)

\item[{$A[i,j]$}]
the  element of the matrix $A$ at $i$-th row, $j$-th column\\
矩陣 $A$ 的第 $i$ 列，第 $j$ 行的元素

\end{SymEntry}


\newpage
%% 論文本體頁碼回復為阿拉伯數字計頁，並從頭起算
\pagenumbering{arabic}
%%%%%%%%%%%%%%%%%%%%%%%%%%%%%%%%